% Generated draft: requires mathematical and source review.
\documentclass[11pt]{article}
\usepackage[margin=25mm]{geometry}
\usepackage{fontspec}
\setmainfont{Latin Modern Roman}
\usepackage{amsmath,amssymb,mathtools,unicode-math}
\setmathfont{Latin Modern Math}
\usepackage{enumitem,needspace,xurl,hyperref}
\setlength{\parindent}{0pt}
\setlength{\parskip}{5pt}
\setlength{\emergencystretch}{3em}
\newcommand{\N}{\mathbb{N}}
\newcommand{\Z}{\mathbb{Z}}
\newcommand{\Q}{\mathbb{Q}}
\newcommand{\R}{\mathbb{R}}
\newcommand{\C}{\mathbb{C}}
\newcommand{\F}{\mathbb{F}}
\newcommand{\A}{\mathbb{A}}
\newcommand{\PP}{\mathbb{P}}
\newcommand{\E}{\mathbb{E}}
\newcommand{\eps}{\varepsilon}
\newcommand{\dd}{\,\mathrm{d}}
\begin{document}
\section*{UnsolvedMath ID: 4300003}
\textbf{Source problem code:} catalogue entry
\subsection{Definitions and mathematical statement}
An algebraic \(\Z^d\)-action in a customary compact-group framework is a homomorphism
\[
\alpha:\Z^d\longrightarrow\operatorname{Aut}(X)
\]
for a compact topological group \(X\), where the automorphisms are continuous; equivalently, its \(d\) generating automorphisms commute. A compact abelian group is the usual more specific setting, but the extracted statement does not state whether this restriction is intended. The spaces singled out are compact spaces locally homeomorphic to \(U\times C\), where \(U\) is a manifold chart and \(C\) is a Cantor set, and compact totally disconnected spaces.

Ward's Problem C asks for an analogue of Pesin theory for suitably defined smooth maps on these spaces. In a manifold-times-Cantor model, a possible differentiable structure is leafwise differentiability in the \(U\)-directions, but the source gives no compatible atlas, derivative, regularity class, invariant measure, or integrability assumptions. On a totally disconnected space it does not specify any replacement for tangent directions. The requested analogue could involve Lyapunov exponents or stable and unstable objects, but no precise existence or comparison theorem for them is supplied; none is invented here.

The second question asks whether any algebraic action with \(d>1\) admits meaningful perturbation. An actual deformation \(s\mapsto\alpha_s\) would have to specify the permitted class and topology and maintain
\[
\alpha_s(e_i)\alpha_s(e_j)=\alpha_s(e_j)\alpha_s(e_i)
\qquad(1\le i,j\le d)
\]
for every parameter \(s\). The source does not say whether perturbations must remain algebraic, may be leafwise smooth or arbitrary topological actions, or must be nontrivial up to conjugacy. Thus this is an existential research question whose adjective meaningful lacks a fixed mathematical criterion.

The supplied research summary labels the programme open, mentioning fragmentary higher-rank rigidity tools. Its statement that no suitable complete theory exists is a supplied status assessment, not an impossibility theorem excluding all deformations or all possible analogues.
\subsection{Short English statement}
Develop Pesin-type theory on the nonmanifold spaces of algebraic dynamics, and determine whether higher-rank algebraic actions admit meaningful perturbations. The required smooth structure and perturbation category are part of the unresolved formulation.
\subsection{Sources}
\begin{itemize}
\item Local UnsolvedMath registry, permanent ID \(4300003\), exact source code \texttt{catalogue entry}, ``Analogues of Pesin theory.''
\item Thomas Ward, \emph{Six problems in Algebraic Dynamics} (2006), Problem C; \url{https://www.imath.kiev.ua/~skolyada/kevin.pdf}. Recorded access 2026-07-29; pdf-text extraction.
\item Supplied OPEN-TRIAGE research summary mentions fragmentary higher-rank rigidity progress without naming a specific theorem or supplying additional publication references.
\end{itemize}
\end{document}
